\documentclass[a4paper]{article}
\usepackage[margin=1in]{geometry}
\usepackage{xcolor}
\usepackage{amssymb}
\usepackage[parfill]{parskip}
\usepackage{graphicx}
\usepackage{hyperref}
\usepackage{enumitem}
\usepackage[backend=biber]{biblatex-chicago}
\addbibresource{schmidt.bib}

\usepackage{fbb}

%\pagestyle{empty}

\newcommand{\entry}[6]{
    \begin{minipage}[t]{.05\textwidth}
        \centering #1
    \end{minipage}
    ---
    \begin{minipage}[t]{.05\textwidth}
        \centering #2
    \end{minipage}
    \hfill
    \begin{minipage}[t]{.85\textwidth}
        \textbf{#3}, #4\\
        #5\\
        \vspace{-1em}
        #6
    \end{minipage}
}

\begin{document}

\begin{minipage}[t]{.5\textwidth}
    \huge \sc Dan Andrei Iliescu
\end{minipage}
\hfill
\begin{minipage}[t]{.5\textwidth}
    \flushright
    E-mail: \href{dai24@cam.ac.uk}{\texttt{dai24@cam.ac.uk}}\\
    Website: \href{https://dan-andrei-iliescu.github.io/}{\texttt{dan-andrei-iliescu.github.io}}\\
    Google Scholar: \href{https://scholar.google.com/citations?user=VimcbYgAAAAJ}{\texttt{user=VimcbYgAAAAJ}}\\
    LinkedIn: \href{https://www.linkedin.com/in/dan-andrei-iliescu/}{\texttt{in/dan-andrei-iliescu}}
\end{minipage}

\vspace{2em}

I am a 3rd-year PhD candidate in Machine Learning studying generative models for high-dimensional data that can learn causal representations of the world.

\vspace{1em}

\textsc{\Large Education}

\rule{\textwidth}{1pt}

\entry{Oct 2019}{June 2023}{University of Cambridge}{UK}{PhD in Machine Learning}{
    \begin{itemize}[leftmargin=*]
        \item Full scholarship from the Alan Turing Institute, the UK's national data science institute.
        \item Developing foundational AI models for learning from unstructured, high-dimensional data.
    \end{itemize}
}

\entry{Oct 2018}{July 2019}{University of Cambridge}{UK}{MPhil in Advanced Computer Science}{
    \begin{itemize}[leftmargin=*]
        \item Starred distinction (+80\%) in ``Machine Learning'' and ``Language Modelling'' courses.
        \item Developing AI models for dealing with confounding in high-dimensional data. 
    \end{itemize}
}

\entry{Sep 2015}{July 2018}{University of Manchester}{UK}{BSc (Hons) in Computer Science, UK First Class (1:1)}{
    \begin{itemize}[leftmargin=*]
        \item Dissertation on visualising the learned concepts of Generative Adversarial Networks. 
    \end{itemize}
}

\vspace{1em}

\textsc{\Large Work Experience}

\rule{\textwidth}{1pt}

\entry{Oct 2019}{Now}{University of Cambridge}{UK}{Teaching Assistant}{
    \begin{itemize}[leftmargin=*]
        \item Teaching ``Data Science'', ``Machine learning'', and ``Information Theory'' courses.
    \end{itemize}
}

\entry{Mar 2023}{Now}{ION.gov.ro}{Bucharest, Romania}{Research Consultant}{
    \begin{itemize}[leftmargin=*]
        \item Developing an AI ChatBot connecting voters with their political representatives.
    \end{itemize}
}

\entry{June 2022}{Feb 2023}{Papercup Technologies Ltd}{London, UK}{Machine Learning Scientist Intern}{
    \begin{itemize}[leftmargin=*]
        \item Authored a research paper on controlling text-to-speech generative models.
    \end{itemize}
}

\entry{May 2021}{Oct 2021}{Yokai.ai}{Paris, France}{Research Consultant}{
    \begin{itemize}[leftmargin=*]
        \item Developed deep learning models for automatically assessing image quality.
    \end{itemize}
}

\entry{June 2018}{Sep 2018}{University of Manchester}{UK}{Research Assistant Intern}{
    \begin{itemize}[leftmargin=*]
        \item Co-authored a research paper on specialisation in neural network ensembles.
    \end{itemize}
}

\entry{June 2017}{Sep 2017}{ARM Ltd}{Manchester, UK}{Software Engineering Intern}{
    \begin{itemize}[leftmargin=*]
        \item Co-authored a whitepaper on optimising ML operations for novel CPU architectures.
    \end{itemize}
}

\vspace{1em}

\newpage

\textsc{\Large Publications}

\rule{\textwidth}{1pt}


\begin{enumerate}[itemsep=1em]
    \item \textbf{Iliescu, Dan Andrei}, Devang Mohan, Tian Huey Teh, and Zack Hodari. Controlling High-Dimensional Data With Sparse Input. Under review at ICML 2023.
    \item Mustafa, Aamir, Aliaksei Mikhailiuk, \textbf{Dan Andrei Iliescu}, Varun Babbar, and Rafał K. Mantiuk.``Training a task-specific image reconstruction loss.'' In Proceedings of the IEEE/CVF winter conference on applications of computer vision, pp. 2319--2328. 2022.
    \item \textbf{Iliescu, Dan Andrei}, Aliaksei Mikhailiuk, Damon Wischik, and Rafal Mantiuk. Disentangling Domain and Content. Preprint 2022.
    \item Webb, Andrew, Charles Reynolds, Wenlin Chen, Henry Reeve, \textbf{Dan Andrei Iliescu}, Mikel Lujan, and Gavin Brown. ``To ensemble or not ensemble: When does end-to-end training fail?.'' In Machine Learning and Knowledge Discovery in Databases: European Conference, ECML PKDD 2020, Ghent, Belgium, September 14–-18, 2020, Proceedings, Part III, pp. 109--123. Springer International Publishing, 2021.
    \item \textbf{Iliescu, Dan Andrei}, and Francesco Petrogalli. ``Arm scalable vector extension and application to machine learning.'' Retrieved October (2018).
\end{enumerate}

\vspace{1em}

\textsc{\Large Selected Projects}

\rule{\textwidth}{1pt}

\entry{2021}{}{AI Mosaic}{Award at St Edmund's College, Cambridge Art Contest}{Python algorithm for assembling images using patches from other images.}{
    \begin{itemize}[leftmargin=*]
        \item Website: \href{https://github.com/Dan-Andrei-Iliescu/ai-painting}{/ai-painting}
    \end{itemize}
}

\entry{2019}{2020}{Machine Learning Reading Club}{Computer Laboratory, Cambridge}{Organized weekly ML paper presentations for PhD students in the Computer Lab.}{
    \begin{itemize}[leftmargin=*]
        \item Website: \href{https://www.cl.cam.ac.uk/research/rainbow/readingclub/}{/readingclub}
    \end{itemize}
}

\entry{2018}{}{Visulising GAN Layers}{Undergraduate Thesis}{Interpreting the concepts learned by a GAN.}{
    \begin{itemize}[leftmargin=*]
        \item Website: \href{https://github.com/Dan-Andrei-Iliescu/visualizing-GAN-layers}{/visualizing-GAN-layers}
    \end{itemize}
}

\vspace{1em}

\textsc{\Large Skills}

\rule{\textwidth}{1pt}

\textbf{ML and Data Science:} Python, Pytorch, Tensorflow, Scipy, Sklearn, Pandas, Plotly 

\textbf{Development:} Git, Docker, Bash, Poetry, DVC

\textbf{Some Experience:} Java, C++

\end{document}